\chapter*{Preface}
\addcontentsline{toc}{chapter}{Preface}
\markboth{Preface}{Preface}

These notes were prepared for Renzo's Fall 2026 course at Colorado
State University. Their organizing premise is that toric geometry is
best learned by moving repeatedly among three languages:
\[
\boxed{
\text{algebra}
\quad\longleftrightarrow\quad
\text{geometry}
\quad\longleftrightarrow\quad
\text{combinatorics}.
}
\]
Characters become monomials, cones become affine charts, rays become
divisors, and lattice points become sections. None of these
translations is merely decorative: each language makes visible
features that are difficult to see in the other two.

The book is divided into two parts.

\textbf{Part~I, Basic Concepts}
(Chapters~0--6), builds the three dictionaries themselves. Chapter~0 is
a panoramic return to \(\PP^2\): it introduces most of the recurring
constructions in one familiar example before the later chapters separate
them into definitions and theorems. Chapters~1--6 then pass from tori and
affine toric varieties to fans, divisors, polytopes, and Cox coordinates.
By the end of Part~I, the reader can translate a basic question among the
algebraic, geometric, and combinatorial languages and choose the one in
which it is easiest to solve.

\textbf{Part~II, (More) Advanced Topics} (Chapters~7--10), puts those
dictionaries to work. Its chapters study canonical and Fano geometry,
moment maps and topology, localization and Riemann--Roch, and finally
variation of quotients and the toric minimal model program. The word
``more'' is intentional: these topics use additional machinery, but the
calculations remain grounded in the same cones, fans, divisors, lattice
points, and homogeneous coordinates developed in Part~I. The division is
therefore a change of emphasis rather than a hard boundary. Part~I builds
the translation tools; Part~II uses them to organize deeper geometry.

The notes are deliberately conversational. Chapter overviews state
the central question; dictionary boxes keep the three languages in
view; surface laboratories make the general theory calculable; and
pause, prediction, and guiding-question boxes invite the reader to
anticipate the next step. Every chapter closes with a summary,
comprehension exercises, a cumulative challenge, and notes and
references. The intended rhythm is to compute first, translate
second, and only then compress the pattern into a theorem.

The notation below records the conventions used throughout. When a
chapter introduces temporary coordinates or labels for a particular
example, those local choices are stated explicitly. The index at the end
of the book collects the main constructions, theorems, and recurring
examples; entries under a broad term often point to more specific
subentries.

\vspace{0.5cm}These notes were developed with substantial expository, editorial, and typesetting assistance from OpenAI’s ChatGPT.

\section*{How this class is going to run}


I don't want this course to be a standard "I lecture you listen" kind of class; at the same time I am wanting to cover quite a bit of material in this one semester that we have, so I have thought about the following shared-responsibility structures, one for the basic concepts, one for the advanced topics.

\subsubsection*{Basic concepts}
For each of the basic concepts chapters, I will first lecture for a few hours, aiming to cover all the important concepts covered in the chapter. For each chapter, the last two hours will be the  responsibility of a  group of students.
These two hours should not be a repetition of my previous lectures - but rather the opportunity to fill in gaps I might have left, or to go into more depth into the material. Some of the things you can do during this time are:
\begin{itemize}
    \item present examples of concepts;
    \item work out significant exercises - both in the comprehension and the challange category;
    \item ask me questions about things are still not clear (but be ready to reason about the questions you ask);
    \item describe to the class modifications to the book you have made (more on this later).
\end{itemize}
But really, you can do anything which is meaningful for the goal of the class, which is understanding the relevant concepts in toric geometry. If you are in doubt if your plan is legit, just check with me.
And the  assignment of the chapters is - drum roll...
\begin{description}
    \item[Chapter 1:] Brendan, Kandin, Vivian;
    \item[Chapter 2:] Ayodeji, Eamon, Kobe, Ryan;
    \item[Chapter 3:] Derek, Makenna,  Ross;
    \item[Chapter 4:] Derek, Eamon, Kobe;
    \item[Chapter 5:] Ayodeji, Brendan, Makenna, Ryan;
    \item[Chapter 6:]  Kandin, Ross, Vivian.
\end{description}
\subsubsection*{Advanced topics}
The four chapters of advanced topics, you will be completely in charge. We will devote one week ( = three hours) to each chapter, and you will be leading the class the whole time. You shouldn't think of just covering all the material in the chapter, but make a selection that makes sense for three hours. Of course, I do encourage you to work with me throughout the semester to prepare for your advanced topics week. Do not leave it last minute. While the assignment of topics to people in the basic concepts part was more or less random, here I am trying to match topics with people's interest. If I am mistaken, and some of you would like to trade topics with someone else, no problem, just let me know.

\begin{description}
    \item[Chapter 7:] Derek, Kandin, Kobe;
    \item[Chapter 8:] Makenna, Ryan, Vivian;
    \item[Chapter 9:] Ayodeji, (Ignacio),   Ross;
    \item[Chapter 10:] Brendan, Eamon, (Jacob).
\end{description}

\subsubsection*{The book}

Over the last month I had quite a lot of conversations with ChatGPT, and came up with this book. It is a bit chatty and not super efficient, but I think it overall does a good job of charting an interesting initial path through all the material covered in Cox-Little-Schenck, with a lot of examples and pictures and emphasizing the conceptual rather than the computational aspects of the theory.

While I did put a bit of effort trying to catch missteps, I make no claim of having produced an error free product, and also - more simply - this is the book that {\bf I} came up with. I would like by the end of the semester for this book to be {\bf our} book. So you are encouraged - especially for the chapters you are responsible for - but not necessarily only for those, to proofread and make modification to the textbook so that it serves the best possible purpose for your learning.

\subsubsection*{AI usage}
AI is here and we can't put our head in the sand about it, but rather we need to learn to use it as a tool for better learning (and not as a substitute for learning).
You are totally welcome to use AI to work on the book. You are totally welcome to have conversations with AI as starting points for your understanding.
But there are two things I ask of you:
\begin{enumerate}
    \item don't let the conversations with AI completely replace conversations with your peers and with me; use your AI conversations as a starting point, but always use another human as the endpoint of your understanding - even if it is just to confirm what you have understdood chatting with your favorite AI;
    \item with the exercises, especially the comprehension exercises, spend a little time on your own (or with human friends) before you turn for help to AI. There is value and good training even in being confused and making unsuccessful attempts at solving problems. If and when you then turn to AI for help, try to have a strategy, and to ask questions that lead towards following that strategy, or understanding why it doesn't work and devising a better one, rather than just asking for a solution.
\end{enumerate}
Why does all of this matter anyway?
\subsubsection*{Evaluation}
Grades - ugh - don't we all hate them? I do. Let me premise that I want to give an A to everyone in this class...but (there is always a but) I do want to make you work for it, with the payoff being that you learn toric geometry. So here is what the evaluation will be based on:
\begin{enumerate}
    \item the two basic concept chapters you are responsible for;
    \item the advanced topic presentation you are responsible for;
    \item one meeting (say 30sh minutes) towards the end of the semester, in which we will have a conversation about some of the comprehension exercises in the book. I don't need you to be able to cleanly solve the random exercise I choose, but I do need you to have an idea for how to tackle it, and to be able to describe it to me. The idea with this format is that throughout the semester you should have a look at each of the comprehension exercises, and get to the point of knowing what they are asking and what the first couple steps towards a solution would be.
\end{enumerate}





